\begin{recipe}{Tarte Tatin met witlof}
\kicker[Bijgerecht,Vegetarisch,Frans,Kerst]{Bijgerecht · vegetarisch · Frans · kerst}
\index[register]{Tarte Tatin met witlof}
\index[register]{Witlof}
\index[register]{Bladerdeeg}
\index[register]{Geitenkaas}

\meta{50 minuten \dvd Voor 4 personen \dvd Oven 200\,°C}

\lettrine{D}{eze} hartige tarte Tatin draait het bittere imago van witlof om: door de
karamel wordt het zoet en mals onder een laag knapperig bladerdeeg. Een makkelijk
bijgerecht, bijvoorbeeld bij het kerstdiner, waar ook kinderen die normaal niet van
witlof houden misschien wel voor overstag gaan.

\blockrule{Ingrediënten}
\begin{ingredients}
  \ing{3}{stronken witlof}\index[register]{Witlof}
  \ing{50 g}{roomboter}
  \ing{50 g}{suiker}
  \ing{4 plakjes}{bladerdeeg}\index[register]{Bladerdeeg}
  \ingb{verse tijm}
  \ingb{wat geitenkaas}\index[register]{Geitenkaas}
\end{ingredients}

\blockrule{Bereiding}
\tip{Geen springvorm? Gebruik een ovenschaal met een bodem van bakpapier, zodat de
karamel niet onder de rand vandaan kan lopen.}
\begin{steps}
  \item Verwarm de oven voor op 200\,°C (boven- en onderwarmte).
  \item Verdeel de boter en suiker over de bodem van een ronde springvorm (ca.\ 20 cm)
        en zet ca.\ 10 tot 12 minuten in de oven, tot de boter en suiker samen een
        lichtbruine karamel vormen.
  \item Snijd de stronken witlof in de lengte doormidden. Leg de helften met de
        snijkant naar boven in de karamel.
  \item Leg de 4 plakjes bladerdeeg over de witlof en druk de randen langs de rand
        van de vorm naar binnen.
  \item Bak 20 tot 25 minuten, tot het bladerdeeg gaar en goudbruin is.
  \item Laat de tarte Tatin 5 minuten rusten, leg dan een groot bord op de vorm en
        draai hem voorzichtig om.
  \item Kruimel de tijm en de geitenkaas erover en serveer direct.
\end{steps}
\end{recipe}
